\title{Parameter-Efficient Orthogonal Finetuning via Factorization}
Further mathematical details are elaborated in Appendix~\ref{app:sn_property}. Viewing the task through the lens of information transmission allows for diverse and flexible designs of sparse matrix factorization schemes that confine the parameter budget while still maintaining the capability to produce expressive dense matrices. Additionally, it is yet unclear whether the butterfly structure constitutes the optimal mechanism for information propagation. Our main contributions are as follows:

\begin{itemize}[leftmargin=*,nosep]
    \item We reframe parameter-efficient orthogonal finetuning by developing a novel information transmission framework, which models the construction of dense orthogonal matrices with minimal parameters as a problem of information flow on a grid-like graph structure. Compared to LoRA, orthogonal finetuning is generally superior in maintaining spectral characteristics, as it exactly preserves the spectral norm of the original pretrained weight matrix $\bm{W}^0$. &\bm{B}_2(\frac{d}{2})
    \end{bmatrix}= \tilde{\bm{B}}(d,d)\tilde{\bm{B}}(d,\frac{d}{2})\cdots\tilde{\bm{B}}(d,2),
\end{aligned}
\end{equation}

Here, $\bm{B}_1(\frac{d}{2})$ and $\bm{B}_2(\frac{d}{2})$ denote two butterfly matrices of dimension $\frac{d}{2}$. The practical effectiveness and simplicity of this low-rank reparameterization approach have contributed to its broad adoption~\cite{dettmers2023qlora,zi2023delta,chavan2023one}. For each curve, the same color represents a BOFT configuration with identical block size, while the leftmost point captures the error for BOFT($1$,$b$) (\emph{i.e.}, the conventional block-diagonal OFT with block size $b$). Concretely, building a dense matrix $\bm{R}\in\mathbb{R}^{d\times d}$ through a sequence of $m$ square matrices, i.e., $\thickmuskip=2mu \medmuskip=2mu \bm{R}=\bm{B}_m\bm{B}_{m-1}\cdots\bm{B}_1$, can be interpreted as moving information across a grid consisting of $d\times(m+1)$ nodes (see Figure~\ref{fig:info_trans}). With this framework, we can represent an orthogonal matrix using butterfly matrices by enforcing the orthogonality of every multiplicative component $\tilde{\bm{B}}(d,k)$ for each $k$ present in the butterfly matrix $\bm{B}(d)$. Learning fast linear transformations with butterfly parameterizations has been explored in~\cite{dao2019learning,dao2019kaleidoscope}. \begin{theorem}[Expressivity of BOFT]\label{thm:exp_boft}
    BOFT possesses greater expressivity than OFT with the same block size. Examining the pattern of non-zero entries in the constructed orthogonal matrix underscores the utility of the information transmission perspective for OFT, as we are concerned solely with the locations of trainable non-zero elements in $\bm{R}$, rather than their numeric values. By employing the Cayley transform or equivalently using $2$-dimensional rotations for each block, as established in \cite{liu2021orthogonal,qiu2023controlling}, one can directly parameterize $\tilde{\bm{B}}(d,2)$ to be orthogonal. Nonetheless, it remains an open question how effectively our orthogonal butterfly matrix can approximate an arbitrary orthogonal matrix. \textbf{Identity Initialization in BOFT}. Specifically, if an entry $R_{ij}$ in matrix $\bm{R}$ is zero, it signifies that the $j$-th source node does not transmit to the $i$-th target node. \section{An Information Transmission View on Orthogonal Finetuning}

We begin by reviewing foundational aspects of OFT. Through this method, a $d$-dimensional dense matrix can be expressed as the product of $\mathcal{O}(\log d)$ sparse matrices, each with only $\mathcal{O}(d)$ nonzero elements. In order to develop a comprehensive and accessible framework for analyzing sparsity in orthogonal matrix factorization, we recast the generation of a dense orthogonal matrix within OFT as an information transmission process. To investigate this, we perform a simulation in which we aim to approximate a randomly generated dense orthogonal matrix~\cite{anderson1987generation} of size $512\times 512$. The proposed BOFT method addresses OFT’s parameter overhead by employing butterfly factorization, thereby making it feasible to leverage orthogonal finetuning for a broader array of model adaptation tasks. Furthermore, the utility of OFT beyond its original use case of controlling text-to-image diffusion models~\cite{qiu2023controlling} has yet to be established. Despite its empirical utility, segmenting a neuron's dimensions by index into groups is not conceptually justified, suggesting that a dense orthogonal parameterization may be preferable. By ensuring each factor is orthogonal, this method achieves a highly efficient orthogonal parameterization requiring only $\mathcal{O}(d\log d)$ parameters—substantially fewer than the standard approach. OFT outperforms LoRA with respect to generalization and produces more reliable training dynamics. This exploits the block structure of $\bm{R}(m, b)$, as its butterfly components can be decomposed or permuted into block-diagonal orthogonal matrices of size $2b \times 2b$. In LoRA~\cite{hulora2022}, Dropout is deployed for its low-rank adaptation parameters to regularize learning and mitigate overfitting, a design choice carried over from conventional Dropout techniques~\cite{srivastava2014dropout}. Within this framework, we observe that the butterfly structure adeptly fulfills the goals of sparse orthogonal matrix decomposition. However, the framework’s reliance on high-dimensional orthogonal matrices leads to a large number of trainable parameters. Of these, model finetuning is particularly noteworthy for its conceptual simplicity, robust effectiveness, and, crucially, its lack of added inference delay. Figure~\ref{fig:info_trans} illustrates an example where matrix factorization leads to a feasible configuration using $10$ edges, which is fewer than the full dense orthogonal alternative. This property is evident from the singular value decomposition $\bm{W}^0 = \bm{U} \bm{\Sigma} \bm{V}^\top$, where $\bm{U}$ and $\bm{V}$ are orthogonal matrices and $\bm{\Sigma}$ is diagonal containing the singular values. Drawing on theoretical insights relating hyperspherical energy to generalization~\cite{liu2018learning,liu2021orthogonal}, orthogonal finetuning~\cite{qiu2023controlling} was introduced, which instead learns an orthogonal transformation for neurons within a layer. For enhanced parameter efficiency, a block-diagonal form is imposed, expressing $\bm{R}$ as $\text{diag}(\bm{R}_1,\bm{R}_2,\cdots,\bm{R}_r)$, with each block $\bm{R}_i\in\mathcal{R}^{b\times b},\forall i$, and satisfying $br=d$. 
\begin{abstract}
Large foundation models are becoming increasingly prevalent, yet the high cost of training them from the ground up presents significant obstacles. More broadly, BOFT($m$, $b$) entails $\frac{1}{2}(b-1)dm$ parameters that can be tuned for adapting a linear layer of dimensions $d \times n$. Expanding this approach, we follow \cite{chen2022pixelated} and define the block butterfly component $\tilde{\bm{B}}^b(d,k)$ such that each element of $\bm{d}_i$ becomes a $b\times b$ matrix. To illustrate, the arrangement in Figure~\ref{fig:info_trans} utilizes $10$ connections, yet the butterfly network accomplishes the same with only $8$. We thus introduce a multiplicative Dropout strategy specific to BOFT: for each forward pass, a random selection is made of $p_1$ percent of the butterfly matrices, and within each chosen butterfly component, $p_2$ percent of the block diagonals, both of which are then replaced with identity matrices. Additionally, the inherent sparsity present in the BOFT matrix factorization may further introduce subtle inductive biases~\cite{gunasekar2017implicit,li2018algorithmic,liu2019neural}. This perspective frames BOFT as learning a similarity matrix $\bm{R}$ under the stringent requirement that $\bm{R}$ remains orthogonal, thereby linking the approach to the broader domains of distance metric learning~\cite{xing2002distance} and bilinear forms~\cite{roman2005advanced}. In contrast, BOFT introduces a fundamentally different finetuning methodology, wherein transformation of the weights is performed via layer-shared weight matrices. Our reasoning builds on the butterfly structure’s flexibility to represent a broad spectrum of classic linear transforms, a capability that the block-diagonal construction in OFT lacks entirely. Addressing the challenge of reducing BOFT’s training time remains an open area for research. Notably, \cite{dao2022monarch} addresses the finetuning of pretrained weights by first mapping them onto a modified version of butterfly matrices, subsequently optimizing these projected components for specific downstream tasks. Overall, BOFT achieves superior parameter efficiency compared to OFT. These edges are mandatory for orthogonality, are non-trainable, and should be excluded from the count of trainable edges since, for instance, in a $d\times d$ orthogonal matrix with $d$ non-zero elements, these entries are restricted to $\pm 1$. Importantly, orthogonality imposes a subtle requirement on the connections between consecutive layers. If the block size is $b = d$, then BOFT($1$, $d$) recovers the full, unconstrained OFT~\cite{qiu2023controlling}. To illustrate, consider the block-diagonal case $\tilde{\bm{B}}(d,2)$ where the block size is $2$. Importantly, upon completion of finetuning, the orthogonal matrices derived through BOFT can be absorbed into the original model without introducing inference overhead. LoRA’s parameter efficiency is achieved by imposing a low-rank constraint, whereas standard OFT attains efficiency by utilizing a block-diagonal structure for the orthogonal matrix~\cite{qiu2023controlling} (with greater efficiency for smaller block sizes). To ensure that the block butterfly component $\tilde{\bm{B}}^b(d,2)$ is orthogonal, each $2b\times 2b$ block is itself parameterized to be orthogonal. A straightforward approach to densely connecting two levels requires $\mathcal{O}(d^2)$ edges (i.e., corresponding to a single dense orthogonal matrix), resulting in $d^2-d$ total connections (for instance, when $d=4$, there are 12 edges). Based on this insight, we introduce Orthogonal Butterfly~(BOFT), a novel finetuning strategy that generalizes OFT within a parameter-efficient and orthogonal context. To overcome this limitation, we analyze OFT through the lens of information transmission and pinpoint several essential criteria for improving parameter efficiency. \textbf{Spectral properties}. Predominant strategies for this purpose include: \emph{model finetuning}~\cite{hulora2022,zaken2022bitfit,guo2020parameter,raffel2020exploring,qiu2023controlling,zhang2022adaptive,chavan2023one}, which preserves the model’s structure while selectively updating a portion of its parameters; \emph{adapter tuning}~\cite{rebuffi2017learning,houlsby2019parameter,pfeiffer2020adapterfusion,lin2020exploring,he2021towards}, which augments the model with extra trainable components; and \emph{prompt tuning}~\cite{li2021prefix,lester2021power}, which attaches learned prefix tokens to model inputs. Distinct from prior works, our approach aims to leverage butterfly structures specifically to improve the parameter efficiency of OFT. Their utility further extends to efficient matrix-vector products~\cite{michielssen1996multilevel,o2010algorithm}, data-sparse matrix representations~\cite{li2015butterfly}, and facilitating network transmission~\cite{klasing1994broadcasting,li2003linear,rajasingh2016transmission}. OFT ensures zero initialization by setting $\bm{R}$ to the identity matrix at the outset. In comparison, standard OFT with a full orthogonal matrix incurs a parameter cost of $\mathcal{O}(d^2)$, and the block-diagonal OFT (with $r$ blocks) requires $\mathcal{O}(bd)$. The butterfly structure has also gained traction as a computer network topology~\cite{solihin2015fundamentals,li2011linear}, offering a solution for streamlined data exchange. Figure~\ref{fig:para_eff} shows the averaged results over 10 different seeds. It is instructive to compare our method with approaches like LoRA~\cite{hulora2022}, which exploit low-rank parameterizations; both low-rank and sparse matrices constitute key classes of structured matrices~\cite{candes2011robust} that drive efficiency in parameter usage. From the information transmission standpoint, two primary criteria emerge: \emph{(i)} \textbf{dense connectivity}, meaning each initial node maintains at least one connection to every final node, and \emph{(ii)} \textbf{minimal free edges}, referring to keeping the total edge count as low as possible within orthogonality constraints. Traditionally, a dense orthogonal matrix of dimension $d$ necessitates $\mathcal{O}(d^2)$ parameters, making direct construction seemingly impractical. In what follows, we describe how leveraging the butterfly structure leads to gains in parameter efficiency. OFT has shown notable advantages in terms of generalizability and convergence, particularly in the context of finetuning text-to-image diffusion models~\cite{qiu2023controlling}. Our objective is to generate a dense orthogonal matrix by assembling $m$ sparse orthogonal matrices, all while minimizing the count of effective, trainable parameters. To allow the butterfly matrix to approximate every orthogonal matrix of size $d$, one can construct products such as $\bm{B}_{d-1,1}(d)\bm{B}^\top_{d-1,2}(d)\cdots\bm{B}_{1,1}(d)\bm{B}^\top_{1,2}(d)$, where all $\bm{B}_{i,j}(d)$ represent butterfly matrices. Because orthogonal matrices require fewer parameters compared to dense matrices, imposing orthogonality introduces further dependencies among edges. While the block-diagonal structure reduces the number of trainable parameters, it imposes a limitation: neuron dimensions (i.e., columns of $\bm{W}^0$) are arbitrarily partitioned into $r$ distinct groups, each independently transformed by a separate orthogonal matrix. For instance, BOFT($9$,$2$) achieves better expressive power than BOFT($1$,$16$), despite requiring significantly fewer parameters. Unlike LoRA, which employs an additive low-rank matrix to update weights, OFT introduces updates through the multiplication of an orthogonal matrix. Specifically, for each $\bm{B}_i$ to achieve full rank—a prerequisite for orthogonality—a set of $d$ connections must exist, guaranteeing a one-to-one mapping between nodes in the $i$-th and $(i=1)$-th layers. \input{rebuttal-results}

\section{Concluding Remarks and Limitations}

This work presents Orthogonal Butterfly, a broadly applicable parameter-efficient finetuning approach for foundation models that leverages the butterfly architecture. This pattern highlights that the butterfly structure forms a more intricate subset within the orthogonal group relative to the block-diagonal approach, making BOFT inherently more expressive than OFT with an equivalent block size. The central concept for enhanced parameter efficiency involves expressing a dense orthogonal matrix as the product of several sparse orthogonal matrices. Thus, ongoing advancements in the field are predicated on the capacity to adapt pre-existing foundation models without full retraining. Thus, achieving an optimal balance between expressivity and regularization is crucial for successful model finetuning. BOFT extends the range of finetuning parameters via structural inductive biases, helping to navigate the trade-off between flexibility and regularity more effectively. More generally, for a dense matrix $\bm{R}\in\mathbb{R}^{d\times d}$ to result from such a product, the set of $m$ matrices $\bm{B}_m,\cdots,\bm{B}_1$ must define a network of directed edges within a $d\times (m+1)$ grid, where each edge only bridges adjacent levels (columns), guaranteeing that any source node on the first level has a path to any node on the $(m+1)$-th level. The construction draws inspiration from the butterfly graphs~\cite{cooley1965algorithm} central to the Cooley-Tukey algorithm, which facilitate dense interconnection between $d$ sources and $d$ destinations with $\mathcal{O}(d\log d)$ edges. We support our method through comprehensive empirical evaluations involving the adaptation of large vision transformers, expansive language models, and text-to-image diffusion systems across a diverse range of vision and language downstream tasks. In this work, we examine Orthogonal Finetuning (OFT), a theoretically grounded finetuning methodology for task adaptation. The orthogonal matrix $\bm{R}(m, b)\in\mathbb{R}^{d\times d}$ is constructed through the product of several orthogonal butterfly matrices. Taking $d=2^N$, we introduce the $d$-dimensional \emph{butterfly matrix} $\bm{B}(d)\in\mathbb{R}^{d\times d}$ recursively as

\begin{equation}
\begin{aligned}
    \!\!\bm{B}(d)=\tilde{\bm{B}}(d,d)\!\cdot\!\begin{bmatrix}
    \bm{B}_1(\frac{d}{2})\!\!\!\!\! Figure~\ref{fig:comp_lora} offers a visual contrast between the two approaches. For illustration, Figure~\ref{fig:info_trans} examines the factorization $\bm{R}=\bm{B}_5\bm{B}_4\bm{B}_3\bm{B}_2\bm{B}_1$, displaying the corresponding sparsity structures and the graph generated by the multiplication. Such butterfly matrices have been used to parametrize orthogonal matrices, enhancing Gaussian elimination by removing the need for pivoting and boosting computational efficiency~\cite{parker1995random}; they have also been leveraged to improve the stability of recurrent neural network training~\cite{jing2017tunable} and in the context of kernel approximation~\cite{munkhoeva2018quadrature}. The LoRA approach~\cite{hulora2022} modifies pretrained weights by incorporating a low-rank matrix product, leading to strong results on NLP benchmarks. This approach allows for an analysis of the parameter-efficiency of OFT from a graph-theoretical viewpoint, making it straightforward to devise feasible factorizations. Such invariance in the spectral norm has been empirically linked to improved training robustness and enhanced model generalization~\cite{miyato2018spectral,yoshida2017spectral}. This construction allows for a computationally efficient parameterization of orthogonal matrices composed of many $2\times 2$ orthogonal submatrices. As foundation models (\emph{e.g.}, \cite{brown2020language,radford2021learning,rombach2022high,kirillov2023segment}) continue to grow in both size and capability, developing methods for their efficient adaptation to new tasks with minimal parameter updates has emerged as an active research area~\cite{treviso2023efficient,ding2023parameter,lialin2023scaling}. As shown in Figure~\ref{fig:blk_diag}, the block-diagonal nature of $\bm{R}$ in standard OFT results in parameter savings, yet inherently limits $\bm{R}$ from being dense. By integrating this design with OFT, we develop Orthogonal Butterfly~(BOFT), a new finetuning approach that optimizes parameter efficiency. The choice of $\bm{R}(m,b)$ within BOFT typically corresponds to a specifically structured subset within the orthogonal group, restricting the overall set of hypotheses the model may express. This exponential growth poses significant barriers for researchers aiming to train such models independently. Numerous studies~\cite{dao2019learning,dao2019kaleidoscope,chen2022pixelated,dao2022monarch} have explored sparse training through butterfly parameterizations. To further contextualize the matrix factorization process, we reinterpret the synthesis of a dense orthogonal matrix as a problem of information dissemination, specifically modeled as information passing across a grid-like graph structure. The butterfly architecture, when combined with appropriate permutations, is capable of exactly reconstructing various well-known fast linear transforms~\cite{dao2019learning,dao2019kaleidoscope} (\emph{e.g.}, fast Fourier transform, Hadamard transform). Common practices in model finetuning, for example, involve using a low learning rate, an empirical strategy intended to keep alterations between the original and adapted models minimal. The information transmission framework provides valuable insights into the sparsity structure of the factorization matrices within OFT. \end{theorem}

Theorem~\ref{thm:exp_boft} implies a straightforward pathway to generalize BOFT: the ultimate orthogonal matrix can be defined as $\thickmuskip=2mu \medmuskip=2mu \bm{R}^G(m_1,b_1,m_2,b_2,l)=\bm{R}_{l,1}(m_1,b_1)\bm{R}_{l,2}^T(m_2,b_2)\cdots\bm{R}_{1,1}(m_1,b_1)\bm{R}_{1,2}^T(m_2,b_2)$, with $\bm{R}_{i,j}^T(m,b)$ standing for the orthogonal matrices employed within BOFT. Consequently, effective adaptation strategies for these sophisticated models to specific downstream tasks have become paramount. \textbf{Inductive bias and generalization in BOFT}. This figure unravels the composite matrix into a network structure: here, each $\bm{B}_i$ serves as an adjacency matrix specifying links between nodes at the $i$-th and $(i+1)$-th levels. Consequently, this block-diagonal configuration, while efficient, may introduce potentially problematic inductive biases (such as limited expressiveness or inability to approximate conventional linear operations), and crucially, optimal determination of the sparsity pattern remains an open challenge. The rationale for this information-theoretic viewpoint stems from the idea that a $d$-dimensional dense matrix functions as a fully connected mapping from $d$ input nodes to $d$ output nodes. Both OFT and BOFT apply an orthogonal transformation $\bm{R}$ from the left, leading to updated weights $\bm{R}\bm{U}\bm{\Sigma}\bm{V}^\top$. These investigations typically involve directly reparameterizing weight matrices with butterfly structures and training neural networks from initialization. Thus, expressing the dense matrix $\bm{R}$ as a product $\bm{B}_m\bm{B}_{m-1}\cdots\bm{B}_1$ equates to modeling sequential information transfer, where each $\bm{B}_i$ represents connectivity at a different stage, with transmission progressing from $\bm{B}_1$ up to $\bm{B}_n$. We validate this claim through experiments spanning computer vision and natural language processing, where BOFT consistently surpasses existing state-of-the-art approaches in terms of both parameter efficiency and generalization. This operation leaves the largest singular value—equivalently, the spectral norm—unaltered. This strategy leverages the idea that computational resources and parameter storage can be traded off, allowing for a decrease in the number of trainable parameters. \textbf{Comparison to butterfly-based sparse training}. We introduce the notion of a \emph{butterfly component} as $\thickmuskip=2mu \medmuskip=2mu \tilde{\bm{B}}(d,k)=\text{diag}(\bm{B}^F_1(k),\cdots,\bm{B}^F_{d/k}(k))$, which is a block-diagonal matrix of size $d\times d$, with each block having size $k$. Model finetuning fundamentally relies on maintaining similarity between the pretrained and the finetuned models according to specified criteria, thereby preserving the acquired knowledge from pretraining. Notably, both block-diagonal and butterfly configurations rely on introducing sparsity into the matrix representation, thereby lowering the trainable parameter count. To mitigate this issue, OFT incorporates a block-diagonal design that reduces parameter counts. The original forward computation, $\bm{z}=(\bm{W}^0)^\top\bm{x}$, is thus altered to $\bm{z}=(\bm{R}\bm{W}^0)^\top\bm{x}$, where the input and output vectors are $\bm{x}\in\mathbb{R}^{d}$ and $\bm{z}\in\mathbb{R}^{n}$, respectively. \section{Orthogonal Parameterization by Butterfly Factorization}

Originally, the butterfly architecture was introduced within the Cooley-Tukey algorithm to enable efficient computation of the fast Fourier transform. Because the orthogonal matrix is formed as a product of sparse factors, these matrix multiplications must be executed iteratively at each step of training. Considering the structured composition of weight matrices, alternative methods aim to safeguard the relational properties within these matrices—most notably, the angles between neuron pairs. Our proposed information transmission framework makes it possible to draw insights from the domain of computer networking, where the study of network topologies for effective information transfer is well established. Configurations with smaller $b$ and larger $m$ generally exhibit increased parameter efficiency. The radix-2 Cooley-Tukey algorithm simplifies the $N$-point discrete Fourier transform into a sequence of two $\frac{N}{2}$-point transforms recursively, producing what is known as a butterfly structure—effectively a composition of sparse matrices, collectively referred to as a butterfly matrix. Drawing inspiration from the efficiency of the Cooley-Tukey fast Fourier transform, we introduce a novel orthogonal parameterization framework based on butterfly structures. \end{abstract}

\section{Introduction}

Models such as ChatGPT~\cite{brown2020language,chen2021evaluating} and Stable Diffusion~\cite{rombach2022high} have highlighted the extraordinary generalization capacities of large foundation models. Choosing $m = \log \frac{2d}{b}$ and $b$ recovers a dense orthogonal transformation with $\bm{R}$ as a block butterfly matrix $\bm{B}^{\frac{b}{2}}(d)$. This surge in model performance is accompanied by a dramatic escalation in model scale

(\emph{e.g.}, GPT-3 contains approximately 175 billion parameters). Consequently, BOFT imposes an inductive bias. For BOFT, we initialize all butterfly matrices as identity matrices (meaning the associated skew-symmetric matrices used in their Cayley transforms are zero). \textbf{Multiplicative Dropout in BOFT}. By assembling these elements, the forward operation in BOFT can be expressed as

\begin{equation}
    \begin{aligned}\nonumber
    \bm{z} = \big(\bm{R}(m, b) \cdot \bm{W}^0\big)^\top \bm{x},~~~~\text{s.t.}~\bigg\{ \bm{R}(m, b) = \prod_{i=1}^m \tilde{\bm{B}}^b_{(d, i)}~~\&~~\underbrace{\big(\tilde{\bm{B}}^b_{(d, j)}\big)^\top\tilde{\bm{B}}^b_{(d, j)} = \tilde{\bm{B}}^b_{(d, j)}\big(\tilde{\bm{B}}^b_{(d, j)}\big)^\top = \bm{I}_d}_{\forall j\in [1, m]}\bigg\},
    \end{aligned}
\end{equation}

Here, for clarity, we write $\tilde{\bm{B}}^b(d, 2^{m - i + 1})$ as $\tilde{\bm{B}}^b_{(d, i)}$, where $\bm{I}_d$ represents the $d \times d$ identity matrix. To maintain the orthogonality constraint during fine-tuning, we adopt Cayley parameterization as described in \cite{qiu2023controlling,liu2021orthogonal}, wherein $\bm{R}=(\bm{I}+\bm{Q})(\bm{I}-\bm{Q})^{-1}$ and $\bm{Q}$ is skew-symmetric, satisfying $\bm{Q}=-\bm{Q}^\top$. In contrast to OFT, which leverages a block-diagonal structure to balance between expressivity and regularization, BOFT utilizes the butterfly structure to provide a continuous and smoother transition between full orthogonal matrices (maximizing expressivity) and identity matrices (emphasizing regularity). In detail, element $(j_1, j_2)$ of $\bm{B}_i$ identifies whether a directed edge runs from node $j_2$ on level $i$ to node $j_1$ on level $(i+1)$ (an entry of zero denotes no connection). In contrast, a non-zero $R_{ij}$ entry permits information flow between those nodes. This mandates at least $d$ inter-layer edges (see Figure~\ref{fig:info_trans} for an example with 4). The butterfly network induces a specialized form of sparse matrix decomposition known as butterfly factorization. \item For the first time, we demonstrate the viability of orthogonal finetuning across a variety of tasks beyond controllable text-to-image generation~\cite{qiu2023controlling}, thereby establishing BOFT as a promising and generalizable model finetuning method. \section{Intriguing Insights and Discussions}

\textbf{Expressivity of BOFT}. Numerous PEFT strategies have been introduced~\cite{houlsby2019parameter,aghajanyan2020intrinsic,hulora2022,edalati2022krona,wang2022adamix,gheini2021cross,zaken2022bitfit,guo2020parameter,sung2021training,ansell2022composable,lester2021power,li2021prefix,vu2022spot,he2021towards,mao2021unipelt,karimi2021compacter,liu2022few,sung2022lst,chen2022parameter,jia2022visual,chen2022adaptformer,zhang2022neural,jie2023fact,lian2022scaling,luo2023towards}, with reparameterization-based techniques~\cite{aghajanyan2020intrinsic,hulora2022,edalati2022krona} being most pertinent to the scope of this work. Thus, orthogonality can occasionally lead to additions or subtractions in feasible edge sets. For example, BOFT($6$,$4$) matches the approximation error of BOFT($2$,$16$) while utilizing substantially fewer parameters. However, this advantage comes with a higher count of trainable parameters than LoRA, suggesting the need for more compact variants. Nevertheless, BOFT is not devoid of shortcomings. Building upon this premise, the Orthogonal Finetuning~(OFT)~\cite{qiu2023controlling} framework was developed, wherein all neurons within a given layer undergo a transformation via a shared orthogonal matrix, which ensures that pairwise neuron angles are rigorously retained throughout the adaptation process. To facilitate the identification of suitable matrix decompositions, we introduce a graphical framework grounded in information transmission principles. Although OFT achieves strong generalization, it still requires substantial trainable parameters, primarily due to the large dimensions associated with orthogonal matrices. Choosing $m_1=m_2=\log d$, $b_1=b_2=2$, and $l=d-1$, the matrix $\bm{R}^G(m_1,b_1,m_2,b_2,l)$ is expressive enough to cover the entirety of the orthogonal group. This matrix composition is sometimes referred to as the kaleidoscope hierarchy~\cite{dao2019kaleidoscope}. To illustrate, consider a $2\times 2$ orthogonal matrix: if the $(1,1)$ entry is zero, the orthogonality constraint forces the $(2,2)$ entry to be zero as well (\emph{i.e.}, omitting one edge enforces omission of another). Here, the y-axis corresponds to the approximation error, and the x-axis indicates the number of effective trainable parameters. To realize parameter reduction within this information dissemination framework, we are inspired by butterfly networks commonly employed in the Cooley-Tukey fast Fourier transform algorithm~\cite{cooley1965algorithm}, where efficient data exchange between graph nodes is a hallmark~\cite{klasing1994broadcasting}. For succinctness, we refer to BOFT configured with $\bm{R}(m, \frac{b}{2})$ as BOFT($m$, $b$), given that $b \geq 2$. BOFT can also be interpreted as learning a bilinear similarity function of the form $\bm{w}^0_i\bm{R}\bm{x}$, where $\bm{w}^0_i$ represents the $i$-th column (neuron) of the weight matrix $\bm{W}^0$. A central challenge lies in constructing dense orthogonal matrices in a manner that remains efficient in terms of parameter count. Setting $m = 1$ yields BOFT($1$, $b$), which coincides with the block-diagonal OFT~\cite{qiu2023controlling} where the block size is $b$. As such, generating a dense orthogonal matrix within the standard OFT necessitates $b = d$, but BOFT provides flexibility to achieve this for arbitrary $b$. We contend that the inductive bias originating from butterfly factorization assists generalization, as its structure is reminiscent of well-established linear transforms, such as the discrete Fourier, sine/cosine, and Hadamard transforms~\cite{dao2019kaleidoscope}. We anticipate that future work could identify even more efficient network configurations for constructing dense orthogonal matrices. To ensure $\bm{B}_5\bm{B}_4\bm{B}_3\bm{B}_2\bm{B}_1$ is fully dense, every initial node must be able to pass information to every node in the sixth layer. The underlying matrix factorization additionally allows for notable weight interpolation properties (see Figure~\ref{fig:boft_interpolation}). & \bm{0} \\
    \bm{0}\!\!\!\!\! \item Drawing inspiration from butterfly structures, notably as deployed in the Cooley-Tukey algorithm, we introduce Orthogonal Butterfly—a new orthogonal finetuning approach that is both parameter-efficient and grounded in the information transmission framework. The results corroborate the overall effectiveness of BOFT as a universal method for model finetuning. Additionally, both~\cite{chen2022pixelated,dao2022monarch} demonstrate that butterfly matrices allow for highly efficient neural network training. Efficient adaptation methods are vital for leveraging these models in various downstream applications. However, due to the multiplicative manner in which BOFT updates weights, standard Dropout does not directly transfer. Since LoRA imposes a fixed rank for all of its inserted matrices, subsequent works~\cite{zhang2022adaptive,valipour2022dylora,zhang2023increlora} have introduced strategies to adaptively tune these ranks per layer to optimize parameter allocation. BOFT serves as a broad generalization of OFT, encompassing it as a particular instance. This mapping between graph connectivity and sparsity applies as well for further reduced products like $\bm{R}=\bm{B}_3\bm{B}_2\bm{B}_1$ or $\bm{R}=\bm{B}_2\bm{B}_1$. Here, each $\bm{B}^F_i(k)$ corresponds to the butterfly factors outlined in Equation~\ref{eq:but_factor}. Nevertheless, we point out that enhanced expressivity does not automatically translate to improved finetuning outcomes, as demonstrated by full finetuning, which is fully expressive yet frequently exhibits suboptimal empirical performance. The main rationale for using orthogonal transformations in finetuning is their ability to maintain pairwise angular relationships among neurons~\cite{liu2018learning,liu2021orthogonal,qiu2023controlling}, thereby preserving a substantial portion of the pretrained model’s semantic knowledge. We circumvent this limitation by innovatively assembling a dense orthogonal matrix through the composition of multiple sparse matrices in a factorized manner. This structural choice narrows the effective hypothesis space within the orthogonal group for subsequent learning tasks. In the context of Fourier transforms, a localized modification in the frequency domain may affect the entirety of the spatial domain—a phenomenon closely related to our goal in information transfer, where each node in the initial layer can relay information to every node in the final layer. Since its resulting orthogonal matrix is formed by the product of multiple orthogonal factors, there is a slight increase in training runtime compared to OFT. \end{itemize}

\section{Related Work}

\textbf{Parameter-efficient finetuning (PEFT)}. \textbf{Orthogonal finetuning as learning bilinear similarity}. For components $\tilde{\bm{B}}^b(d,k)$ with $k>2$, their non-zero patterns correspond to block-wise permutations of those in $\tilde{\bm{B}}^b(d,2)$, enabling similar orthogonal parameterizations. Despite its parameter savings, this structure imposes a rigid sparsity pattern on the orthogonal matrix and applies the orthogonal transformations independently within each block. Since the pattern of non-zero entries in any butterfly component results from permuting the non-zero pattern of $\tilde{\bm{B}}(d,2)$, each butterfly component can be analogously made into an orthogonal matrix. Specifically, OFT optimizes an orthogonal matrix $\bm{R}\in\mathbb{R}^{d\times d}$ for a pretrained linear mapping $\bm{W}^0\in\mathbb{R}^{d\times n}$. Owing to the butterfly structure’s established role in efficient linear transformations—including the discrete Fourier and discrete cosine transforms—we posit that this approach to parameter efficiency naturally embeds an implicit inductive bias, enhancing both generalization and resilience to overfitting. For butterfly matrices, specifically $m = \log d$, $b = 2$, BOFT’s parameter count is $\mathcal{O}(d\log d)$. Extensive experiments highlight the empirical advantages of BOFT in adapting large language models, substantial vision models, and text-to-image generative models. This leads to a key research question: \emph{Is it possible to devise a dense orthogonal matrix structure that maintains parameter efficiency?}

To investigate this issue, we suggest constructing a dense orthogonal matrix by multiplying several sparse orthogonal matrices. When focusing on a partial factorization, such as $\bm{R}=\bm{B}_4\bm{B}_3\bm{B}_2\bm{B}_1$, which concerns nodes in levels one and five, we observe that node 1 at the source does not reach node 3 at the target, which is reflected by a zero at the $(3,1)$ entry in the associated sparsity pattern. In OFT, the fine-tuning process involves reparameterizing the weight matrix by expressing it as the product of a trainable orthogonal matrix and the static pretrained weight matrix. Most parameter-efficient adaptation approaches maintain initialization aligned with the pretrained network so that adaptation is stable and non-destructive; for instance, LoRA initializes its low-rank matrices to zero. Given $k \geq 2$ as a power of $2$, we formally define the \emph{butterfly factor} $\bm{B}^F(k)$ by

\begin{equation}\label{eq:but_factor}
\begin{aligned}
    \bm{B}^F(k)=\begin{bmatrix}
    \text{diag}(\bm{d}_1) & \text{diag}(\bm{d}_2) \\
    \text{diag}(\bm{d}_3) & \text{diag}(\bm{d}_4)
    \end{bmatrix}\in\mathbb{R}^{k\times k},
\end{aligned}
\end{equation}

where each vector $\bm{d}_i\in\mathbb{R}^{\frac{k}{2}},\forall i$ specifies diagonal elements. \textbf{Butterfly structures}. \item We deliver several theoretical arguments to explain how BOFT achieves a substantial reduction in trainable parameter count without sacrificing expressive capacity or training stability.